%------------------------------------------------------------------------------%

\usepackage{calculo_varias_variaveis-1}

\title{Propriedades do Vetor Gradiente}

%------------------------------------------------------------------------------%
\begin{document}

\addtitle

%------------------------------------------------------------------------------%
\section{Vetor Gradiente}

%------------------------------------------------------------------------------%
\begin{frame}{Vetor Gradiente}
  Se uma função \(f(x, y)\) possui derivadas parciais,
  \pause
  seu \emph{vetor gradiente} é
  \[
    \pause
    \nabla f
    \pause
    \;=\; \D{f}{x}\veci + \D{f}{y}\vecj
    \pause
    \;=\; \left(\begin{array}{c} \D{f}{x} \\[5mm] \D{f}{y} \end{array}\right)
  \]
\end{frame}

%------------------------------------------------------------------------------%
\begin{frame}{Vetor Gradiente 3D}
  Se uma função \(f(x, y, z)\) possui derivadas parciais,
  \pause
  seu \emph{vetor gradiente} é
  \[
    \pause
    \nabla f
    \pause
    \;=\; \D{f}{x}\veci + \D{f}{y}\vecj + \D{f}{z}\veck
    \pause
    \;=\; \left(\begin{array}{c} \D{f}{x} \\[5mm] \D{f}{y} \\[5mm] \D{f}{z} \end{array}\right)
  \]
\end{frame}

%------------------------------------------------------------------------------%
\section{Propriedades}

%------------------------------------------------------------------------------%
\begin{frame}{Propriedades da Derivada Direcional}
  Podemos usar a fórmula do produto interno
  \[
    \left(D_u f\right)
    \pause
    = \nabla f \cdot u
    \pause
    = \abs{\nabla f} \abs{u} \cos(\theta)
    \pause
    = \abs{\nabla f} \cos(\theta)
  \]
\end{frame}

%------------------------------------------------------------------------------%
\begin{frame}{Propriedades da Derivada Direcional}
\begin{enumerate}[<+->]
  \setlength\itemsep{1.5em}
  \item A função aumenta mais rapidamente quando $\cos(\theta)=1$, $\theta=0$ \\
        Direção do gradiente
  \item A função descresce mais rapidamente quando $\cos(\theta)=-1$, $\theta=\pi$ \\
        Direção oposta ao gradiente
  \item Se $u$ for ortogonal a \(\nabla f\neq 0\) ele aponta para uma direção de variação zero
\end{enumerate}
\end{frame}

%------------------------------------------------------------------------------%
\addtocounter{exemplo}{1}
\begin{frame}{Exemplo \theexemplo}
  Encontre as direções nas quais \(f(x, y) = x^2 + \frac{y^2}{2}\)
  \begin{enumerate}
    \item[a)] cresce mais rapidamente no ponto \((1, 1)\)
    \item[b)] descresce mais rapidamente no ponto \((1, 1)\)
    \item[c)] quais as direções de variação zero de $f$ em \((1, 1)\)
  \end{enumerate}
\end{frame}

\begin{frame}{Exemplo \theexemplo}
  Calculando o gradiente
  \[
    \nabla f
    \pause
    = \D{f}{x} \veci + \D{f}{y}\vecj
    \pause
    = \D{}{x}\left(x^2 + \frac{y^2}{2}\right) \veci
    + \D{}{y}\left(x^2 + \frac{y^2}{2}\right) \vecj
    \pause
    = 2x\veci + y\vecj
    \pause
    = \vetor{2x}{y}
  \]
  \pause
  Avaliando o gradiente no ponto solicitado
  \[
    \pause
    \nabla f (1, 1)
    \pause
    = \left(2x\veci + y\vecj\right)\evalat{(1, 1)}{}
    \pause
    = \vetor{2}{1}
  \]
\end{frame}

\begin{frame}{Exemplo \theexemplo}
  \begin{enumerate}
  \item[a)]
    \pause
    cresce mais rapidamente no ponto \((1, 1)\)
    \[
      \pause
      v_M
      \pause
      = \nabla f(1, 1)
      \pause
      = \vetor{2}{1}
    \]
  \item[b)]
    \pause
    descresce mais rapidamente no ponto \((1, 1)\)
    \[
      \pause
      v_m
      \pause
      = -\nabla f (1, 1)
      \pause
      = \vetor{-2}{-1}
    \]
\end{enumerate}
\end{frame}

\begin{frame}{Exemplo \theexemplo}
\begin{enumerate}
  \item[c)]
    \onslide<+->{quais as direções de variação zero de $f$ em \((1, 1)\)} \\
    \onslide<+->{$v_z$ ortogonal ao gradiente}
    \begin{align*}
      \onslide<+->{\nabla f(1, 1) \cdot n & = 0 \\[2mm]}
      \onslide<+->{(2\veci + \vecj) \cdot (n_1\veci + n_2\vecj)   & = 0 \\[2mm]}
      \onslide<+->{2n_1 + n_2                                     & = 0 \\[2mm]}
      \onslide<+->{n_2                                            & = -2n_1}
    \end{align*}
    \onslide<+->{As direções são}
    \[
      \onslide<+->{n = \vetor{1}{-2}}
      \onslide<+->{
      \qquad\text{e}\qquad
      -n = \vetor{-1}{2}}
    \]
\end{enumerate}
\end{frame}

%------------------------------------------------------------------------------%
%\addtocounter{exemplo}{1}
%\begin{frame}{Exemplo \theexemplo}
%\end{frame}
%
%\begin{frame}{Exemplo \theexemplo}
%\end{frame}

%------------------------------------------------------------------------------%
\section{Gradientes e Curvas de Nível}

%------------------------------------------------------------------------------%
\begin{frame}{Curvas de Nível}
  Curva de Nível é o conjunto de pontos \((x, y)\) onde
  \[
    f(x, y) = c
  \]
  para um valor constante $c$

  \pause
  \vspace{0.5\baselineskip}
  Em alguns casos, é possível escrever essa região como uma curva paramétrica
  \[
    x = g(t)
    \qquad
    y = h(t)
  \]
  \pause
  Vetorialmente temos
  \[
    r(t) = g(t)\veci + h(t)\vecj
  \]
\end{frame}

%------------------------------------------------------------------------------%
\begin{frame}{Tangente de uma Curva de Nível}
  \begin{columns}
  \column{0.5\linewidth}
    \begin{align*}
      \onslide<+->{f\big(g(t), h(t)\big) & = c \\[2mm]}
      \onslide<+->{\frac{d}{dt} f\big(g(t), h(t)\big) & = 0 \\[2mm]}
      \onslide<+->{\D{f}{x}\frac{dg}{dt} + \D{f}{y}\frac{dh}{dt} & = 0  \\[2mm]}
      \onslide<+->{\left(\D{f}{x}\veci + \D{f}{y}\vecj\right) \cdot
      \left(\frac{dg}{dt}\veci +\frac{dh}{dt} \vecj\right) & = 0 \\[2mm]}
      \onslide<+->{\nabla f \cdot \frac{dr}{dt} & = 0}
    \end{align*}
  \column{0.5\linewidth}
    \onslide<+->{O gradiente é ortogonal ao vetor tangente à curva de nível}
  \end{columns}
\end{frame}

%------------------------------------------------------------------------------%
\begin{frame}{Funções Diferenciáveis}
  Em todo ponto \((x_0, y_0)\) no domínio de uma \emph{função diferenciável}
  \(f(x, y)\),

  \pause
  o \emph{gradiente} de $f$
  \pause
  é normal à \emph{curva de nível} que passa por \((x_0, y_0)\)
\end{frame}

\framedgraphic{Gradientes e Curvas de Nível}{gradiente_curva_nivel}{}

%------------------------------------------------------------------------------%
\addtocounter{exemplo}{1}
\begin{frame}[t]{Exemplo \theexemplo}
  Encontre uma equação para a tangente à eclipse
  \[
    \frac{x^2}{4} + y^2 = 2
  \]
  no ponto \((-2, 1)\)
\end{frame}

\begin{frame}[t]{Exemplo \theexemplo}
  Encontre uma equação para a tangente à eclipse
  \[
  \frac{x^2}{4} + y^2 = 2
  \]
  no ponto \((-2, 1)\)

  \addgraphic{9}{6.5,2.5}{tangente_elipse}
\end{frame}

\begin{frame}{Exemplo \theexemplo}
  Precisamos notar que a elipse é uma curva de nível da função
  \[
    f(x, y) = \frac{x^2}{4} + y^2
  \]

  \pause
  Calculando as derivadas parciais de $f$ temos
  \[
    \pause
    \D{f}{x}
    \pause
    = \D{}{x}\left(\frac{x^2}{4} + y^2\right)
    \pause
    = \frac{x}{2}
  \]
  \[
    \pause
    \D{f}{y}
    \pause
    = \D{}{y}\left(\frac{x^2}{4} + y^2\right)
    \pause
    = 2y
  \]
\end{frame}

\begin{frame}{Exemplo \theexemplo}
  Então o gradiente de $f$, no ponto \((-2, 1)\), é
  \pause
  \[
    \left(\nabla f\right) \evalat{(-2, 1)}{}
    \pause
    = \left(\frac{x}{2}\veci + 2y\vecj\right) \evalat{(-2, 1)}{}
    \pause
    = -\veci + 2\vecj
  \]
\end{frame}

\begin{frame}{Exemplo \theexemplo}
  \onslide<+->{A reta tangente é ortogonal ao vetor gradiente}
  \begin{align*}
    \onslide<+->{\left(\nabla f\right) \evalat{(x_0, y_0)}{}
                 \cdot \big((x-x_0)\veci + (y-y_0)\vecj\big) & = 0 \\[3mm]}
    \onslide<+->{\left(\nabla f\right) \evalat{(-2, 1)}{}
                 \cdot \big((x+2)\veci + (y-1)\vecj\big) & = 0 \\[3mm]}
    \onslide<+->{\left(-\veci + 2\vecj\right) \cdot \big((x+2)\veci + (y-1)\vecj\big) & = 0 \\}
    \onslide<+->{-(x+2) + 2(y-1) & = 0 \\}
    \onslide<+->{x +2y -4 & = 0 \\}
    \onslide<+->{x +2y & = 4}
  \end{align*}
\end{frame}

%------------------------------------------------------------------------------%
\begin{frame}{Propriedades Algebricas do Gradiente}
  \begin{enumerate}[<+->]
    \setlength\itemsep{1.5em}
    \item Soma                        \tabto{55mm} \(\nabla(f+g) = \nabla f + \nabla g\)
    \item Diferença                   \tabto{55mm} \(\nabla(f-g) = \nabla f - \nabla g\)
    \item Multiplicação por constante \tabto{55mm} \(\nabla(kf) = k\nabla f\)
    \item Produto                     \tabto{55mm} \(\nabla(fg) = f\nabla g + g\nabla f\)
    \item Quociente                   \tabto{55mm}
      \(\nabla\left(\frac{f}{g}\right) = \frac{g\nabla f - f\nabla g}{g^2}\)
  \end{enumerate}
\end{frame}

%------------------------------------------------------------------------------%
\addtocounter{exemplo}{1}
\begin{frame}{Exemplo \theexemplo}
  Dadas \(f(x, y)=x^2-2y\) e \(g(x, y)=xy^2\)
  \begin{enumerate}
    \item calcule os gradientes de $f$ e $g$
    \item calcule \(\nabla(f-g)\)
    \item calcule \(\nabla(fg)\)
  \end{enumerate}
\end{frame}

\begin{frame}{Exemplo \theexemplo}
  \[
    \nabla f
    \pause
    \;=\; \D{f}{x}\veci + \D{f}{y}\vecj
    \pause
    \;=\; 2x\veci - 2\vecj
  \]

  \pause
  \vspace{\baselineskip}
  \[
    \nabla g
    \pause
    \;=\; \D{g}{x}\veci + \D{g}{y}\vecj
    \pause
    \;=\; y^2\veci + 2xy\vecj
  \]
\end{frame}

\begin{frame}{Exemplo \theexemplo}
  \begin{align*}
    \onslide<+->{\nabla(f-g)}
    \onslide<+->{& = \nabla f - \nabla g \\[2mm]}
    \onslide<+->{& = \big(2x\veci - 2\vecj\big) - \big(y^2\veci + 2xy\vecj\big) \\[2mm]}
    \onslide<+->{& = \big(2x-y^2\big)\veci + \big(-2-2xy\big)\vecj  \\[2mm]}
    \onslide<+->{& = \big(2x-y^2\big)\veci - 2\big(1+xy\big)\vecj \\[2mm]}
    \onslide<+->{& = \vetor{2x-y^2}{-2-2xy}}
  \end{align*}
\end{frame}

\begin{frame}{Exemplo \theexemplo}
  \begin{align*}
    \onslide<+->{\nabla(fg)}
    \onslide<+->{& = f\nabla g + g\nabla f \\[2mm]}
    \onslide<+->{& = \big(x^2-2y\big)\big(y^2\veci + 2xy\vecj\big)
      + \big(xy^2\big)\big(2x\veci + 2\vecj\big) \\[2mm]}
    \onslide<+->{& = \big(x^2-2y\big)\big(y^2\veci\big)
      + \big(x^2-2y\big)\big(2xy\vecj\big)
      + \big(xy^2\big)\big(2x\veci\big)
      + \big(xy^2\big)\big(2\vecj\big) \\[2mm]}
    \onslide<+->{& = \big(x^2y^2-2y^3\big)\veci
      + \big(2x^2y^2\big)\veci
      + \big(2x^3y-4xy^2\big)\vecj
      + \big(2xy^2\big)\vecj \\[2mm]}
    \onslide<+->{& = \big(3x^2y^2 - 2y^3\big)\veci
      + \big(2x^3y - 2xy^2\big)\vecj \\[2mm]}
    \onslide<+->{& = \vetor{3x^2y^2 - 2y^3}{2x^3y - 2xy^2}}
  \end{align*}
\end{frame}

%------------------------------------------------------------------------------%
\section{Lista Mínima}

%------------------------------------------------------------------------------%
\begin{frame}{Lista Mínima}
  Cálculo Vol. 2 do Thomas 12\textsuperscript{a} ed. -- Seção 14.5
  \begin{enumerate}
    \item Estudar o texto da seção
    \item Resolver os exercícios: 2, 4, 6, 8, 10, 12, 14, 16, 18, 20, 22, 24, 26, 28, 29
  \end{enumerate}

  \vfill
  Atenção: A prova é baseada no livro, não nas apresentações
\end{frame}

%------------------------------------------------------------------------------%
\end{document}
%------------------------------------------------------------------------------%
